\documentclass[11pt,a4paper,fontset=windows]{ctexart}
% ===== 公共导言区：高等代数（一）期中模拟卷 =====
\usepackage{amsmath,amssymb}
\usepackage{unicode-math}
\setmathfont{TeX Gyre Termes Math}
\usepackage[margin=2.05cm,top=2.35cm,bottom=2.0cm,headsep=11pt,includefoot]{geometry}
\usepackage{fancyhdr}
\usepackage{lastpage}
\usepackage{enumitem}
\usepackage{array}
\usepackage{bm}
\usepackage{xcolor}
\usepackage{zhnumber}
\usepackage{needspace}

\definecolor{hdgray}{gray}{0.32}
\definecolor{lngray}{gray}{0.5}
\definecolor{notebg}{gray}{0.955}

\pagestyle{fancy}
\fancyhf{}
\renewcommand{\headrulewidth}{0.4pt}
\renewcommand{\footrulewidth}{0pt}
\fancyhead[L]{\small\color{hdgray}\songti 高等代数（一）\ 期中模拟试卷（\thepp）}
\fancyhead[R]{\small\color{hdgray}\songti 考试时间：120 分钟\quad 满分：100 分}
\fancyfoot[C]{\small\color{lngray}\songti 第 \thepage\ 页\quad 共 \pageref{LastPage}\ 页}

\setlength{\parindent}{0pt}
\setlength{\parskip}{2pt}
\linespread{1.16}

% 题目列表 1—8 连续编号
\newlist{qs}{enumerate}{1}
\setlist[qs]{label=\textbf{\arabic*.},ref=\arabic*,leftmargin=2.55em,
  labelsep=0.55em,itemsep=23pt,topsep=6pt,parsep=0pt}
\newlist{note}{enumerate}{1}
\setlist[note]{label=(\arabic*),leftmargin=1.75em,itemsep=1.5pt,topsep=2.5pt,parsep=0pt}
\newlist{sub}{enumerate}{1}
\setlist[sub]{label=(\arabic*),leftmargin=1.9em,itemsep=3pt,topsep=3pt,parsep=0pt}

\newcommand{\secA}[1]{\Needspace{9\baselineskip}\par\vspace{10pt}{\large\bfseries\songti #1}\par\vspace{3pt}\nopagebreak}
\newcommand{\blank}[1][2.6cm]{\underline{\hspace{#1}}}
\newcounter{pp}
\renewcommand{\thepp}{\chinese{pp}}

\setcounter{pp}{5}
\begin{document}

\begin{center}
{\LARGE\bfseries\songti 高等代数（一）期中模拟试卷（五）}\\[7pt]
{\normalsize\songti 适用范围：第一章\ 多项式（\S1\,--\,\S9）；第二章\ 行列式（\S2.1\,--\,\S2.7）}\\[3pt]
{\small\color{hdgray}\songti 多元多项式、对称多项式不在考查范围}
\end{center}
\vspace{-4pt}

\begin{center}
\begin{tabular}{@{}p{0.335\textwidth}p{0.335\textwidth}p{0.26\textwidth}@{}}
\songti 姓名：\blank[2.5cm] & \songti 学号：\blank[2.5cm] & \songti 得分：\blank[1.9cm]
\end{tabular}
\end{center}

\vspace{2pt}
\hrule height 0.5pt
\begin{center}
\setlength{\fboxsep}{6pt}%
\colorbox{notebg}{\parbox{0.94\textwidth}{\small\songti
\textbf{考生须知：}
\begin{note}
\item 本试卷共 8 题，满分 100 分，考试时间 120 分钟。
\item 计算题须写出必要的计算过程，证明题须给出完整的推理步骤；只写答案不给分。
\item 所有答案一律写在答题纸上，写在试卷上无效。
\end{note}}}
\end{center}
\hrule height 0.5pt
\vspace{4pt}

\secA{一、计算题（每小题 10 分，共 40 分）}
\begin{qs}
\item 设
\[
f(x)=\begin{vmatrix}
x+1&2&3\\
1&x+2&3\\
1&2&x+3
\end{vmatrix}.
\]
求 $f(x)$ 在 $\mathbb{C}$ 中的全部根，并写出 $f(x)$ 在 $\mathbb{R}[x]$ 中的标准分解式。

\item 设 $n\ge 3$。计算 $n$ 阶行列式
\[
D_n=\begin{vmatrix}
1&2&3&\cdots&n\\
2&3&4&\cdots&n+1\\
3&4&5&\cdots&n+2\\
\vdots&\vdots&\vdots&&\vdots\\
n&n+1&n+2&\cdots&2n-1
\end{vmatrix}.
\]

\item 用克拉默法则解下列线性方程组：
\[
\begin{cases}
x_1+2x_2+3x_3=9,\\[2pt]
3x_1-x_2+x_3=10,\\[2pt]
x_1-3x_2+2x_3=11.
\end{cases}
\]

\item 设 $a,b$ 为实参数，$f(x)=x^{4}+ax^{2}+b$。
\begin{sub}
\item 求使 $f(x)$ 有重根的全部 $(a,b)$；
\item 对 $a^{2}=4b$ 且 $a<0$ 的情形，写出 $f(x)$ 的重根。
\end{sub}
\end{qs}

\secA{二、证明题（每小题 15 分，共 60 分）}
\begin{qs}[resume*]
\item 设 $f(x)\in P[x]$，$\deg f=n\ge 1$。证明：
\[
f(x)\mid f'(x)\iff f(x)=a(x-b)^{n},
\]
其中 $a\in P$，$a\neq 0$，$b\in P$。

\item 设 $f(x)\in\mathbb{R}[x]$，且对任意实数 $x$ 都有 $f(x)\ge 0$。证明：存在
$u(x),v(x)\in\mathbb{R}[x]$，使得
\[
f(x)=u^{2}(x)+v^{2}(x).
\]

\item 设 $f(x)\in P[x]$，$\deg f=n\ge 0$。
\begin{sub}
\item 证明：$f(x)$ 在 $P$ 中至多有 $n$ 个根（重根按重数计算）；
\item 由此证明：若 $f(x)$ 在无穷多个 $x\in P$ 处取值为零，则 $f(x)\equiv 0$。
\end{sub}

\item 设 $A$ 为 $n$ 阶整数矩阵（元素全为整数），且 $|A|=\pm 1$。证明：$A^{-1}$ 也是整数矩阵。
\end{qs}

\end{document}
